% Source: https://artofproblemsolving.com/community/c3h1786975_mathcounts_problem
% Stable identity: aops-topic-1786975-q01
Whether a number lies in $K$ depends only on its remainder modulo $28=\operatorname{lcm}(4,7)$. Among $1,\dots,28$ there are $7$ multiples of $4$, $4$ multiples of $7$ and $1$ multiple of both, so by inclusion--exclusion $28-7-4+1=18$ numbers lie in $K$. Each block of $28$ consecutive integers therefore contains $18$ members of $K$.

Since $79=4\cdot18+7$, the first $4$ blocks ($1$ to $112$) contain $72$ members, and we need the $7$th member after $112$. Skipping $116$, $119$ and $120$, the members are
\[
113,\ 114,\ 115,\ 117,\ 118,\ 121,\ 122,\ \dots
\]
so the $79$th member is $\boxed{122}$.
